\section{Campo de Velocidade do Som}
\label{sec:campo_de_velocidade_do_som}

Embora seja comum considerar a velocidade do som na água como uma constante, em torno de $1500 \si{\meter/\second}$, essa aproximação é inválida para ambientes de grande escala \cite{kieffer_sound_1977,trusler_determination_2017,li_sound_2015,pinson_sound_2010}. Muitas propriedades podem afetar a velocidade do som na água, como temperatura, pressão estática, e salinidade, além de correntes marítimas \cite{wilson_speed_1959,medwin_fundamentals_1999,mcdougall_getting_2011,delgrosso_new_1974}. À medida que essas propriedades variam na massa de água, a velocidade também varia, criando um \gls{ssf} na massa de água. Embora o termo \gls{ssp} \cite{li_acoustic_2019,li_sound_2015,pinson_sound_2010} também seja comumente usado, ele geralmente se refere ao cenário de estimativa 1D.

%\subsection{Propagação sonora na água}
\subsection{Princípios}

O conhecimento da velocidade do som em cada ponto de uma massa de água permite o conhecimento das propriedades --- como caminho, fase, amplitude, e tempo de trânsito --- de uma onda que se propaga por esse meio.

\subsubsection*{Relação do SSF com parâmetros TPS}

Como já citado, as propriedades de \gls{tps} afetam a velocidade do som na água, atuando de maneiras diferentes. Essa relação é conhecida, e trabalhos iniciais que tratavam da reconstrução do \gls{ssf} visavam estimar essas propriedades principais, bem como a relação entre elas e a velocidade do som. Trabalhos como \cite{wilson_speed_1959,medwin_fundamentals_1999,mcdougall_getting_2011,delgrosso_new_1974,chen_speed_1977,wong_speed_1995,roquet_accurate_2015,mackenzie_nineterm_1981,wilson_equation_1960} operaram dentro dessa estrutura, alguns ainda sendo usados atualmente para modelar o comportamento de correntes marítimas em escala global.

Modelos como os desenvolvidos por \cite{wilson_equation_1960,mackenzie_nineterm_1981,delgrosso_new_1974} relacionam diretamente as variáveis de \gls{tps} com a velocidade do som por meio de funções polinomiais, tendo sido obtidas empiricamente com base em expedições de medição dos parâmetros \gls{tps}, e da velocidade do som. Estes modelos são considerações formulações fechadas, pois a \gls{ssf} pode ser calculada usando um número finito de operações padrão.

O modelo TEOS-10 \cite{mcdougall_getting_2011} é outro método para estimação do \gls{ssf}, resolvendo equações químico-termodinâmicas para determinar a distribuição de velocidade adequada, em função da energia. Na modelagem são resolvidas uma série de equações interconectadas que envolvem a composição química detalhada da água e o equilíbrio térmico. Também baseia-se em parâmetros de \gls{tps} modificados (temperatura conservativa, pressão hidrostática e salinidade absoluta), mas não é uma formulação fechada (como nos modelos desenvolvidos por \cite{wilson_equation_1960,mackenzie_nineterm_1981,delgrosso_new_1974}).

\subsubsection*{Refração}

Um efeito importante, quando tratando de ondas sonoras que atravessam um ambiente com velocidade sonora variável, é o efeito da refração, em que o caminho que a onda percorre não é necessariamente retilíneo. Por estarmos trabalhando com corpos d'água de grandes dimensões (podendo chegar à casa de quilômetros de largura e/ou profundidade), a velocidade do som também pode variar muito e de forma imprevisível ao longo do corpo d'água, necessitando que a refração seja considerada.

Esse efeito está fortemente ligado ao princípio de Fermat de tempo mínimo \cite{se-yuen_closer_1986}, que dita que uma onda (modelada por um raio) que atravessa um meio sempre buscará o menor tempo de viagem. Isso faz com que conhecer o SSF torna o tempo de primeira chegada da onda previsível. Se assumirmos que o \gls{ssf} varie discretamente de velocidade (uma aproximação que será explorada adiante), o princípio de Fermat pode ser representado pela lei de Snell \cite{perrine_fermats_1961}, descrevendo a mudança de ângulo (refração) da onda ao atingir uma interface (discreta) entre dois meios com velocidade do som diferentes.

Por outro lado, a refração pode ser tratada pelo princípio de Huygens, \cite{bormann_seismic_2012}, que modela cada ponto em uma frente de onda como uma fonte sonora secundária. A interação entre essas fontes sonoras virtuais produz o envelope da onda ao longo do tempo e, com isso, gera os mesmos efeitos de refração.

\subsection{Estimação do SSF por inversão}
\label{subsec:sec1:estimacao_ssf_por_inversao}

Como discutido, é possível modelar o \gls{ssf} a partir dos parâmetros de \gls{tps}. Contudo, a estimação desses parâmetros é muitas vezes inviável na prática. Dessa forma, a maioria das abordagens modernas para a reconstrução do \gls{ssf} evita estimar essas propriedades e trata a estimação como um problema de inversão. Parte-se do princípio de que, como o conhecimento do \gls{ssf} permite calcular a trajetória e as propriedades de uma onda que atravessa o meio; então o conhecimento dos caminhos (e/ou das propriedades) permite estimar o \gls{ssf}.

A maioria das técnicas utiliza fontes e receptores sonoros dentro do corpo d'água e, comparando as propriedades das ondas produzidas e recebidas, estima o caminho das ondas; e, a partir disso, obtém uma reconstrução para o \gls{ssf}. Diferentes técnicas tratam esse problema de inversão usando uma vasta gama de tipos de dados e medições, bem como abordagens matemáticas e modelagens distintas.

\subsubsection*{Tomografia de Tempo de Trânsito}

A \gls{ttt} \cite{bormann_seismic_2012,stefanov_travel_2019} é um dos métodos de inversão mais simples para estimativa do \gls{ssf}. Ela utiliza a diferença no tempo de chegada dos sinais acústicos nos sensores, exigindo uma distribuição espacial de fontes e receptores para resolução espacial suficiente.

Ela se baseia em uma discretização espacial do \gls{ssf} em células (ou às vezes camadas), em que a velocidade do som é considerada constante dentro de cada célula; e no princípio de Fermat de tempo mínimo, relacionando o \gls{ssf} ao tempo de primeira chegada da frente de onda em cada ponto do espaço.

No processo iterativo que rege a \gls{ttt}, calcula-se o comprimento dos percurso das ondas em cada célula, resultando em um sistema de equações usado para estimar/iterar o \gls{ssf}. Adaptações do \gls{ttt} incluem a hipótese de o \gls{ssf} se comportar aproximadamente como conhecido, isso sendo inserido como uma regularização, parametrização, ou permitindo uma decomposição modal.

Suas principais vantagens são robustez, simplicidade, e eficiência computacional, a último permitindo implementação em tempo real. No entanto, a \gls{ttt} depende de um posicionamento e sincronização exata dos dispositivos (fontes e receptores), e sua suposição de discretização resulta em resolução limitada.

\subsubsection*{\ingles{Matched Field Inversion}}

Na técnica \gls{mfi} \cite{dosso_bayesian_2011}, é utilizada (primariamente) a informação de pressão sonora medida nos receptores. Essa informação é comparada com modelos de pressão acústica para diferentes \gls[SSFs]{ssfs}, e esses \gls[SSFs]{ssf} são adaptados e iterados de forma a minimizar o erro entre os dados de pressão sonora obtidos e a pressão sonora estimada.

Essas técnicas são semelhantes à formação de feixe (\textsl{beamforming}), utilizada em diversas áreas como acústica e processamento de sinais \cite{schmidt_environmentally_1990,brienzo_broadband_1993}, sendo que ambas --- \gls{mfi} e \textsl{beamforming} --- exploram a correlação espacial entre os receptores e suas medições, modelando o campo acústico dessa modelagem espacial.

O \gls{mfi} também emprega a decomposição modal dos sinais, com cada modo se comportando de maneira diferente na água. Dessa forma, ela incorpora informações mais complexas sobre os sinais, como amplitude e fase, dependendo não apenas do tempo de trânsito. Entre suas vantagens destacam-se o melhor aproveitamento dos dados obtidos, a resolução espacial aprimorada (em comparação com o \gls{ttt}) e a maior robustez ao ruído devido ao processamento coerente dos sinais obtidos nos receptores. No entanto, ela requer informações espaço-temporais sobre a fonte e os receptores, como diretividade e comportamento da potência em função da frequência. Além disso, não é adequada para corpos d'água muito grandes, visto que um corpo maior resultaria em um modelo de pressão sonora com uma quantidade exponencialmente maior de parâmetros. Essa técnica também é computacionalmente muito custosa.

\subsubsection*{\ingles{Full-Waveform Inversion}}

As técnicas de \gls{fwi} \cite{virieux_overview_2009,fichtner_multiscale_2013} visam a reconstrução da forma de onda dos sinais presentes no \gls{ssf}, considerando diversos comportamentos que são negligenciados pelas técnicas mencionadas anteriormente. Nesse método, o erro entre as formas de onda estimadas e medidas é minimizado iterativamente, buscando o \gls{ssf} que melhor represente os dados obtidos. Essa técnica pode operar no domínio do tempo, ou no da frequência.

Dado isso, a técnica \gls{fwi} está entre as que melhor utilizam os dados disponíveis nos receptores, resultando no \gls{ssf} mais fiel, além de aproveitar ao máximo os sinais disponíveis. Entretanto, devido à consideração de interações complexas entre os sinais e múltiplos caminhos de onda, há um aumento significativo na complexidade matemática e computacional. Isso, aliado à não linearidade do problema, resulta em uma técnica muito complexa e por vezes inviável, que também requer uma estimativa inicial adequada dos parâmetros para sua robustez e convergência.

\subsection{Limitações do projeto relevantes à estimação do SSF}

Dentro do presente projeto do OBNZip 2, surgiram algumas restrições quanto aos dispositivos (fontes e receptores) que serão utilizados, e tipo de informação que pode ser obtida.

\begin{enumerate}[itemsep=0pt]
	\item Transceptores são alocadas ao longo do leito oceânico.
	\item Transceptores são alocados próximo à superfície, em profundidades não superiores a $500\si{\meter}$ em situações de água profunda.
	\item Fontes se comunicam acusticamente com receptores, e o tempo de trânsito dessa comunicação é medido.
	\item Apenas a informação de tempo de trânsito é obtida na comunicação (considerando os dispositivos atualmente alocados).
\end{enumerate}

Embora os dispositivos alocados no leito e na superfície são ambos transceptores (isto é, ambos emitem e recebem comunicação), será convencionado que os dispositivos no leito oceânico serão \gls{fontes}, e os na superfície serão \gls{receptores}. Isso simplifica o modelo matemático e a descrição do problema, e deve ser considerado futuramente ao utilizar os desenvolvimentos apresentados.

\subsubsection*{Fontes sonoras}

É desejado utilizar os \gls[OBNs][nós de leito oceânico][Ocean Bottom Nodes]{obn} (inclusive especificamente \ingles{On Demmand} \gls[OBNs]{obn} da Sonardyne) que já estão dispostos no leito oceânico como as fontes sonoras, de forma a não necessitar a construção e implantação de novos equipamentos.% Isso leva a outras restrições, como quanto às frequências sonoras que podem ser utilizadas (limitadas às produzidas pelos OBNs).

\subsubsection*{Receptores}

Os receptores serão implementados por meio de expedições com \ingles{streamers}, \gls[ROVs][veículos operados remotamente][Remotely Operated Vehicles]{rov} ou outros dispositivos, em que esses se comunicam com os \gls[OBNs]{obn}. Embora o problema seja matematicamente modelado como utilizando diversos receptores, na prática isso é implementado como um único (ou alguns poucos) receptor(es), com esse sendo alocado em posições diferentes.

Uma restrição relacionada aos receptores é quanto ao seu posicionamento: dado que o custo de uma expedição está atrelado à profundidade que os receptores serão utilizados (mais caro levá-los até áreas mais profundas, além de uma complexidade maior), deseja-se limitar a profundidade na qual eles podem ser utilizados. Em específico, é de interesse que os receptores estejam limitados a profundidades de $500\si{\meter}$, como alinhado com a Petrobras.

\subsubsection*{Comunicação e tempo de trânsito}

Como o interesse inicial é a utilização da comunicação acústica entre os dispositivos (ao menos nesse primeiro momento), a informação disponível que pode ser utilizada é restrita ao tempo de trânsito entre fontes e receptores.

Tecnicamente, o esquema que será utilizado envolve dispositivos transceptores, operando simultaneamente como fonte e receptor: o dispositivo A emite o sinal; o disp. B recebe, processa, e responde; e o disp. A recebe essa resposta, medindo o tempo de viagem A--B--A. Contudo, por questões de simplificação da modelagem, assume-se que o tempo de viagem A--B é igual ao tempo de viagem B--A, bem como que o disp. B é a fonte, e o disp. A é o receptor.

Também é assumido que não ocorra comunicação entre fonte-fonte (não se sabe o tempo de trânsito entre OBNs), nem entre receptor-receptor; apenas fontes e receptores.

\subsection{Diferenças entre os casos 1D e multi-dimensionais} % Adicionado por sugestão do Luis

Anteriormente, foram apresentados os casos 1D e multi-dimensionais, sem que esses forem profundamente tratados, em particular quanto às suas peculiaridades. Trivialmente, a passagem do caso 1D para o 2D (e do 2D para o 3D, e desse para o 4D) leva a um aumento na complexidade do problema, atrelado tanto ao aumento do número de unidades a serem estimadas no ambiente; mas também um aumento na complexidade das interações entre as ondas que trafegam o ambiente e as células que o compõem, com essas podendo viajar em mais direções, e podendo tomar caminhos diferentes.

Além disso, a alocação de dispositivos está intimamente ligada à estrutura que pode ser estimada. Por exemplo, em um caso com dispositivos alocados em uma única coluna e alinhados verticalmente, não haverá medições do ambiente além dessa coluna (em particular de informações de tempo de primeira chegada), impossibilitando qualquer tipo de estimação que não seja uni-dimensional. Portanto, de maneira geral, a dimensionalidade da estimação está limitada à dimensionalidade da alocação dos dispositivos.

\subsection{\textsl{Mindmap}}

A figura da \cref{fig:mindmap_sec1} mostra um \textsl{mindmap} dos conceitos expostos ao longo da \cref{sec:campo_de_velocidade_do_som}, de forma a dispor graficamente a organização das informações apresentadas. A cor de cada nó indica o conceito-pai dele.

\begin{figure}[!t]
	\centering
	\begin{tikzpicture}[
		mindmap,
		text width=2.5cm,
		align=center,
		every node/.style={concept},
		concept color=gray!50,
		level 1/.append style={
			level distance=5cm,
			minimum size=3cm,
			text width=2.8cm,
			font=\large,
		},
		level 2/.append style={
			level distance=3.5cm,
			minimum size=2cm,
			text width=1.8cm,
            font=\footnotesize
		},
		level 3/.append style={
			level distance=3.5cm,
			minimum size=2cm,
			text width=1.8cm,
            font=\footnotesize
		},
		]
		
		\node (root) [minimum size=4.5cm, text width=4.3cm, line width=6pt, fill=gray!30]{\Large\textbf{\textsl{Sound Speed Field (SSF)}}}
		child[concept color=red!10, grow=30,] {
			node{\textbf{Princípios}}
			child[grow=0] { node{Parâmetros \\ TPS} }
			child[grow=-45] { node (refracao) {Refração} }
		}
		child[concept color=green!10!white!95!black, grow=-30] {
			node{\textbf{Métodos de\\estimação}}
			child[grow=315] { node{FWI e MFI \\(Huygens)} }
			child[grow=270] { node (ttt) {TTT \\ (Fermat)} } 
            child[concept color=red!10] {node at (refracao) {Refração} }
		}
		child[concept color=blue!10, grow=-90] {
			node{\textbf{Limitações\\do projeto}}
			child[grow=-195] { node{Fontes} }
			child[grow=-150] { node{Receptores} }
			child[grow=-105] { node{Comunicação} }
			child[grow=-60] { node{Info. de tempo\\de trânsito}
                child[concept color=green!10!white!95!black] { node at (ttt) {TTT \\ (Fermat)} }
            }
		}
		;	
	\end{tikzpicture}
	\caption{\textsl{Mindmap} da \cref{sec:campo_de_velocidade_do_som}.}
	\label{fig:mindmap_sec1}
\end{figure}
